\section{Introduction}
\label{sec:intro}

Parameter-efficient continual learning (PECL) with LoRA~\citep{hu2022lora} has become an appealing strategy for adapting large pretrained vision models to sequential tasks while updating only a small number of parameters~\citep{Wang_2022_CVPR,Gao_2023_ICCV,Smith2023CODA}.
Among PECL variants, SeqLoRA is particularly memory-efficient because it keeps the pretrained backbone frozen and maintains a single shared low-rank adapter across the entire task sequence.
This design avoids task-wise parameter growth, but it also creates a direct mechanism for catastrophic forgetting.
Since every new task updates the same adapter, the effective adapter update $\Delta W=BA$ can drift away from configurations that supported earlier tasks.
Without replay, task-specific adapter isolation, or an explicit stability constraint, the shared adapter has no structural protection against old-task degradation.

Existing methods address this stability-plasticity tension from two complementary directions.
Sharpness-aware optimization, such as SAM and GAM, improves current-task plasticity by biasing training toward flatter regions of the task loss landscape~\citep{foretSharpnessAwareMinimizationEfficiently2021,zhengGradientAwareminimizationDistributed2021}.
Such flatness has also been linked to improved retention in continual learning~\citep{NEURIPS2020_518a38cc,shiOvercomingCatastrophicForgetting2021}.
In contrast, Fisher-based regularization, including EWC and EWC-LoRA, improves stability by penalizing drift along directions that are important for previous tasks~\citep{EWC2017,zheng2025ewclora}.
A natural response is to combine these two remedies.
However, this combination is not automatically principled.
If a sharpness-aware optimizer is applied directly to a summed task-plus-Fisher objective, the perturbation used to estimate flatness is influenced by both the current-task loss and the old-task Fisher penalty.
As a result, the flatness correction no longer measures current-task sharpness alone.
It becomes entangled with the stability constraint, making the resulting update difficult to interpret and potentially weakening the functional separation between plasticity and stability.

This paper asks a method-design question: how should flatness control and Fisher stabilization be combined in continual LoRA adaptation?
Our answer is that they should not be coupled inside a shared sharpness-aware perturbation closure.
Instead, they should be computed as separate gradient channels with distinct roles.
Current-task flatness should be estimated only from the current-task loss, while old-task stability should be imposed independently through a Fisher-weighted drift penalty in the effective adapter-update space.
This separation is especially important in LoRA, where the trainable factors are only a reparameterization of the functional update.
Regularizing raw factor coordinates can be sensitive to the non-uniqueness of the LoRA factorization, whereas constraining $\Delta W=BA$ directly better aligns the stability penalty with the adapter update that affects the model.

We propose \textbf{Flat-Stable LoRA (FS-LoRA)}, a decoupled optimizer for continual LoRA adaptation.
FS-LoRA has three core design choices.
First, it computes an adapter-subspace first-order sharpness correction using only the current-task loss.
This prevents the old-task Fisher penalty from contaminating the perturbation used for flatness estimation.
Second, it applies Fisher regularization in the effective adapter-update space $\Delta W=BA$, rather than in the raw LoRA factor coordinates.
This makes the stability constraint better aligned with the functional adapter update.
Third, it trace-normalizes the Fisher estimate.
This normalization is necessary because raw Fisher values in low-rank adapter space can be extremely small, which can make the Fisher gradient structurally inactive unless the regularization weight is heavily retuned for each dataset, rank, or task sequence.
With trace normalization, the Fisher weight acts more like a relative stability coefficient rather than a compensation factor for arbitrary Fisher scale.
We also provide a constrained optimization interpretation in which FS-LoRA minimizes current-task loss and first-order sharpness subject to a Fisher drift budget.
This view yields an optional task-boundary dual-ascent controller for adapting the Fisher weight when the stability constraint is violated.

The proposed decoupling is supported by a gradient-geometric diagnostic.
We measure the cosine similarity between the first-order flatness correction and the Fisher gradient throughout training.
Across six fine-grained continual recognition datasets, this cosine remains close to zero, approximately between $-0.03$ and $-0.001$.
This finding suggests that the two components usually capture distinct update directions: the flatness correction responds to current-task curvature, whereas the Fisher gradient constrains old-task sensitive drift.
We use this near-orthogonality as a diagnostic that supports the decoupled design, not as a hard assumption enforced by the optimizer.
FS-LoRA does not perform gradient surgery or projection.
Its key design choice is more basic: avoid mixing flatness and stability before the sharpness correction is computed.

\begin{figure}[t]
\centering
\fbox{\rule[-.5cm]{0cm}{4cm} \rule[-.5cm]{8cm}{0cm}}
\caption{
Cosine distribution between the first-order flatness correction and the Fisher stability gradient across all tasks and six fine-grained datasets.
The values cluster near zero, indicating that the two components behave as distinct gradient channels.
This diagnostic supports the FS-LoRA design choice of decoupling current-task flatness control from past-task Fisher stabilization.
}
\label{fig:cosine_histogram}
\end{figure}

Empirically, FS-LoRA shows that decoupled flatness and Fisher stabilization can improve continual LoRA adaptation, although its effect depends on dataset structure, Fisher activity, and the scale of the flatness correction.
On six fine-grained recognition benchmarks with a frozen ViT-B/16 backbone, FS-LoRA improves final average accuracy over SeqLoRA+GAM on CUB200, Flowers, and ImageNet-R, with the largest gain on Flowers.
The method shows mixed behavior on Cars196, Aircraft, and OxfordPet.
The OxfordPet result exposes a failure mode in which the flatness correction can become too large relative to the clean task gradient, suggesting that adaptive perturbation scaling is an important future direction.
On standard PECL benchmarks, the completed ImageNet-A run shows that an LR-tuned FS-LoRA configuration is within one percentage point of the EWC-LoRA reference, while CIFAR-100, ImageNet-R, and DomainNet runs remain preliminary.
These results support the method design while also identifying the regimes where further robustness improvements are needed.

Our contributions are as follows:
\begin{enumerate}
    \item \textbf{A decoupled flat-stable optimizer for continual LoRA adaptation.}
    We propose FS-LoRA, which separates current-task first-order sharpness control from past-task Fisher stabilization.
    Unlike coupled sharpness-aware Fisher regularization, FS-LoRA computes the sharpness correction only from the current-task loss and injects the Fisher gradient as an independent stability component.

    \item \textbf{Trace-normalized Fisher regularization in effective adapter-update space.}
    We define the Fisher penalty and drift reference in $\Delta W=BA$ rather than in raw LoRA factor coordinates, and introduce trace normalization to prevent Fisher gradients from becoming inactive under low-rank parameterization.

    \item \textbf{A constrained optimization interpretation with optional dual control.}
    We formulate FS-LoRA as a penalty-method approximation to minimizing current-task loss and first-order adapter-subspace sharpness subject to a Fisher drift budget.
    This interpretation yields an optional dual-ascent controller for adapting the Fisher weight when the stability constraint becomes active.

    \item \textbf{Gradient-geometric diagnostics and empirical evaluation.}
    We systematically measure the interaction between the flatness correction and the Fisher gradient, finding near-orthogonality across six fine-grained datasets.
    Experiments show that FS-LoRA improves performance in several settings, reveals failure cases caused by flatness-scale mismatch, and provides a practical diagnostic for designing joint plasticity-stability optimizers.
\end{enumerate}